\documentclass[12pt, a4paper]{article}

%% ── Packages ────────────────────────────────────────────────────────────────
\usepackage[utf8]{inputenc}
\usepackage[T1]{fontenc}
\usepackage[spanish]{babel}
\usepackage{geometry}
\usepackage{graphicx}
\usepackage{booktabs}
\usepackage{array}
\usepackage{longtable}
\usepackage{multirow}
\usepackage{amsmath}
\usepackage{amssymb}
\usepackage{hyperref}
\usepackage{xcolor}
\usepackage{fancyhdr}
\usepackage{titlesec}
\usepackage{caption}
\usepackage{subcaption}
\usepackage{enumitem}
\usepackage{float}
\usepackage{microtype}
\usepackage{lmodern}
\usepackage{parskip}

%% ── Page geometry ────────────────────────────────────────────────────────────
\geometry{
  top    = 2.5cm,
  bottom = 2.5cm,
  left   = 2.8cm,
  right  = 2.8cm
}

%% ── Colors ───────────────────────────────────────────────────────────────────
\definecolor{tec}{RGB}{0, 100, 60}
\definecolor{dark}{RGB}{30, 30, 30}
\definecolor{mid}{RGB}{80, 80, 80}
\definecolor{light}{RGB}{240, 245, 240}

%% ── Hyperref ─────────────────────────────────────────────────────────────────
\hypersetup{
  colorlinks = true,
  linkcolor  = tec,
  citecolor  = tec,
  urlcolor   = tec
}

%% ── Section formatting ───────────────────────────────────────────────────────
\titleformat{\section}
  {\large\bfseries\color{dark}}
  {\thesection~|}{0.5em}{}[\vspace{-0.3em}\color{tec}\rule{\linewidth}{0.4pt}]

\titleformat{\subsection}
  {\normalsize\bfseries\color{dark}}
  {\thesubsection~|}{0.5em}{}

\titleformat{\subsubsection}
  {\normalsize\itshape\color{mid}}
  {\thesubsubsection}{0.5em}{}

%% ── Header / footer ─────────────────────────────────────────────────────────
\pagestyle{fancy}
\fancyhf{}
\fancyhead[L]{\small\color{mid} Análisis Topológico de Datos — PPI Berries (WPUSI01102B)}
\fancyhead[R]{\small\color{mid} Topología de Datos, 6.º Semestre}
\fancyfoot[C]{\small\thepage}
\renewcommand{\headrulewidth}{0.3pt}

%% ── Caption style ────────────────────────────────────────────────────────────
\captionsetup{font=small, labelfont=bf, labelsep=period}

%% ── Math commands ────────────────────────────────────────────────────────────
\newcommand{\mae}{\mathrm{MAE}}
\newcommand{\rmse}{\mathrm{RMSE}}

%% ═══════════════════════════════════════════════════════════════════════════
\begin{document}

%% ── Title page ───────────────────────────────────────────────────────────────
\begin{titlepage}
  \centering
  \vspace*{1.5cm}
  {\Large\bfseries\color{tec}
    Análisis Topológico de Series de Tiempo:\\[0.4em]
    Índice de Precios al Productor de Berries (FRED)\\[0.6em]}
  {\large\color{mid} Comparación de modelos con y sin TDA para pronóstico}
  \vspace{1cm}

  \rule{\linewidth}{0.5pt}
  \vspace{0.5cm}

  {\normalsize
    \begin{tabular}{ll}
      \textbf{Curso}        & Topología de Datos \\[2pt]
      \textbf{Período}      & 6.º Semestre, 2026 \\[2pt]
      \textbf{Dataset}      & \texttt{WPUSI01102B} — FRED, Jun 2008 – Abr 2026 \\[2pt]
      \textbf{Institución}  & Tecnológico de Monterrey, Campus Guadalajara \\
    \end{tabular}
  }
  \vspace{0.5cm}
  \rule{\linewidth}{0.5pt}

  \vspace{0.8cm}
  {\normalsize
    \begin{tabular}{ll}
      \textbf{Nombre}               & \textbf{Matrícula} \\[4pt]
      Josué Tapia Hernández         & A01621056 \\
    \end{tabular}
  }

  \vfill
  {\small Guadalajara, Jalisco, México \quad · \quad Junio 2026}
\end{titlepage}

%% ── Table of contents ────────────────────────────────────────────────────────
\tableofcontents
\newpage

%% ── Abstract ─────────────────────────────────────────────────────────────────
\section*{Resumen ejecutivo}
\addcontentsline{toc}{section}{Resumen ejecutivo}

Este trabajo aplica Análisis Topológico de Datos (TDA) al Índice de Precios al Productor de Berries
publicado por la Reserva Federal de Estados Unidos (\texttt{WPUSI01102B}), una serie mensual de 215
observaciones que abarca junio de 2008 hasta abril de 2026. El objetivo es determinar si las
características topológicas derivadas del embedding de Takens y la homología persistente de Vietoris-Rips
mejoran la capacidad de pronóstico respecto a modelos clásicos sin TDA.

Se construyeron y evaluaron siete clases de modelos: Random Forest (RF), SARIMA, y cuatro variantes
de LSTM con distintos conjuntos de características (sin TDA, con TDA v1, con TDA v2 enriquecida,
y con variables exógenas macroeconómicas más mecanismo de atención). El modelo ganador,
\textbf{LSTM + TDA~v2 + Exógenas + Atención}, alcanzó un MAE de 29.11 puntos PPI en el
conjunto de prueba, frente a 41.19 del mejor modelo sin TDA (RF). La incorporación de
características topológicas redujo el error en un 29~\% y el coeficiente de determinación $R^2$
pasó de valores negativos (RF sin TDA) a 0.54.

\textbf{Palabras clave:} análisis topológico de datos, embedding de Takens, homología persistente,
Vietoris-Rips, LSTM, pronóstico de precios agrícolas, PPI berries.

\newpage

%% ═══════════════════════════════════════════════════════════════════════════
\section{Introducción}

\subsection{Descripción del problema}

Los precios de berries (arándanos, fresas, frambuesas y moras) exhiben una volatilidad pronunciada
producida por la confluencia de tres factores: estacionalidad agrícola (ciclos de cosecha anuales),
choques macroeconómicos (crisis financiera de 2008, pandemia de COVID-19, inflación post-pandemia)
y perturbaciones climáticas que alteran la oferta de forma abrupta. Esta combinación genera una
serie de tiempo con estructura cíclica fuerte pero no estacionaria, lo que dificulta el pronóstico
con métodos convencionales.

El Análisis Topológico de Datos (TDA) ofrece una perspectiva complementaria: mediante el embedding
de Takens, la serie univariada se reconstruye como una nube de puntos en un espacio de alta dimensión,
y la homología persistente cuantifica los ciclos y componentes conexas de esa nube. Estas
características topológicas son invariantes frente a deformaciones suaves de la señal y, por tanto,
capturan patrones de largo plazo (estacionalidad, rupturas estructurales) que los descriptores
estadísticos convencionales podrían omitir.

\subsection{Objetivo}

El objetivo general es comparar el desempeño predictivo de modelos \textit{sin} y \textit{con} TDA
para el pronóstico del índice PPI de berries a un horizonte de un mes hacia adelante. Los objetivos
específicos son:

\begin{enumerate}[leftmargin=1.5em]
  \item Caracterizar topológicamente tres subperiodos históricos (2008-2012, 2013-2019, 2020-2026)
        mediante diagramas de persistencia y entropía topológica.
  \item Construir una familia de modelos de pronóstico que van desde un RF sin TDA hasta un LSTM
        con características topológicas enriquecidas y variables exógenas.
  \item Cuantificar la ganancia marginal atribuible a TDA en cada clase de modelo.
  \item Explorar técnicas avanzadas: función de pérdida Enhanced Peak (EP), cabeza multi-tarea de
        clasificación de extremos, y características de calendario y lag.
\end{enumerate}

\subsection{Datos}

La fuente de datos es el \textbf{FRED Producer Price Index for Berries}, identificado con el código
\texttt{WPUSI01102B}. La serie mide el nivel general de precios que reciben los productores
domésticos de berries, usando como base el año 1982 (índice base = 100). La Tabla~\ref{tab:datos}
resume las principales características del dataset.

\begin{table}[H]
  \centering
  \caption{Características del dataset \texttt{WPUSI01102B}.}
  \label{tab:datos}
  \begin{tabular}{ll}
    \toprule
    \textbf{Atributo}       & \textbf{Valor} \\
    \midrule
    Fuente                  & Federal Reserve Bank of St.~Louis (FRED) \\
    Código                  & \texttt{WPUSI01102B} \\
    Frecuencia              & Mensual \\
    Periodo                 & Junio 2008 – Abril 2026 \\
    Observaciones           & 215 \\
    Columnas                & \texttt{observation\_date}, \texttt{WPUSI01102B} \\
    Unidad                  & Índice PPI (base 1982 = 100) \\
    Media                   & 142.47 \\
    Desviación estándar     & 44.71 \\
    Mínimo                  & 69.2 (Febrero 2020) \\
    Máximo                  & 299.4 (Enero 2021) \\
    Valores faltantes       & 0 \\
    \bottomrule
  \end{tabular}
\end{table}

La serie presenta estacionalidad anual marcada, con picos en los meses de noviembre-diciembre
(temporada de cosecha alta en Chile y México, principal fuente de importaciones invernales) y
valles en los meses de verano en el hemisferio norte. El rango de variación es amplio: el índice
ha oscilado entre 69.2 y 299.4 puntos, reflejando tanto la volatilidad estructural del mercado
como los choques excepcionales de la pandemia de COVID-19.

La Tabla~\ref{tab:subperiodos} desglosa las estadísticas descriptivas por subperiodo histórico,
definidos según los regímenes económicos identificados en el análisis exploratorio.

\begin{table}[H]
  \centering
  \caption{Estadísticas descriptivas por subperiodo histórico.}
  \label{tab:subperiodos}
  \begin{tabular}{lcccccc}
    \toprule
    \textbf{Subperiodo} & \textbf{Fechas} & \textbf{$n$} & \textbf{Media} & \textbf{Desv. Est.} & \textbf{Mín.} & \textbf{Máx.} \\
    \midrule
    Crisis financiera   & Jun 2008 – Dic 2012 & 55 & 148.3 & 47.2 & 80.5  & 282.7 \\
    Estabilidad         & Ene 2013 – Dic 2019 & 84 & 135.3 & 39.1 & 74.2  & 248.5 \\
    Pandemia/Inflación  & Ene 2020 – Abr 2026 & 76 & 146.2 & 47.6 & 69.2  & 299.4 \\
    \midrule
    \textbf{Total}      & Jun 2008 – Abr 2026 & 215 & 142.5 & 44.7 & 69.2 & 299.4 \\
    \bottomrule
  \end{tabular}
\end{table}

%% ═══════════════════════════════════════════════════════════════════════════
\section{Metodología}

\subsection{Análisis Topológico de Datos en series de tiempo}

El TDA aplicado a series de tiempo sigue el paradigma de reconstrucción de atractores propuesto
por Takens (1981). A partir de una serie univariada $\{p_t\}$, se construye un embedding en
$\mathbb{R}^d$ mediante el mapeo:

\begin{equation}
  \Phi: p_t \mapsto (p_t,\, p_{t-\tau},\, p_{t-2\tau},\, \ldots,\, p_{t-(d-1)\tau})
\end{equation}

donde $d$ es la dimensión de embedding y $\tau$ es el retardo temporal. La nube de puntos
resultante $\mathcal{E} \subset \mathbb{R}^d$ preserva las propiedades topológicas del atractor
subyacente.

Sobre $\mathcal{E}$ se construye el complejo simplicial de Vietoris-Rips parametrizado por un radio
$\epsilon$, y se calcula su homología persistente: para cada dimensión $k$ (H$_0$ = componentes
conexas, H$_1$ = ciclos), se rastrean los pares de nacimiento/muerte $(\beta^k,\, \delta^k)$ a
medida que $\epsilon$ crece desde 0 hasta $\infty$. Los pares de larga vida $\delta^k - \beta^k \gg 0$
señalan estructura topológica robusta (no debida a ruido).

La Tabla~\ref{tab:tda_params} resume los parámetros utilizados en este trabajo. Para el análisis
subperiódico (Sesiones 2-4) se empleó $d = 12$ (un año de datos mensuales) y $\tau = 3$
(trimestral). Para el pipeline de pronóstico (Sesiones 6-10) se utilizó una ventana deslizante
de $W = 36$ meses con $d = 6$, $\tau = 3$, produciendo embeddings de 21 puntos por ventana.

\begin{table}[H]
  \centering
  \caption{Parámetros TDA utilizados.}
  \label{tab:tda_params}
  \begin{tabular}{lcc}
    \toprule
    \textbf{Parámetro}         & \textbf{Análisis subperiódico} & \textbf{Pronóstico (rolling)} \\
    \midrule
    Dimensión de embedding $d$ & 12 & 6 \\
    Retardo $\tau$             & 3  & 3 \\
    Ventana $W$                & Subperiodo completo & 36 meses \\
    Puntos por embedding       & Variable & 21 \\
    Dimensiones homológicas    & H$_0$, H$_1$ & H$_0$, H$_1$ (v2: +H$_2$) \\
    Biblioteca                 & \texttt{giotto-tda} & \texttt{giotto-tda} \\
    \bottomrule
  \end{tabular}
\end{table}

El análisis de homología persistente reveló diferencias topológicas significativas entre los
tres subperiodos (Tabla~\ref{tab:persistencia_h1}).

\begin{table}[H]
  \centering
  \caption{Persistencia máxima en H$_1$ por subperiodo (embedding $d=12$, $\tau=3$).}
  \label{tab:persistencia_h1}
  \begin{tabular}{lccc}
    \toprule
    \textbf{Subperiodo} & \textbf{Persistencia máx. H$_1$} & \textbf{Ciclos detectados} & \textbf{Interpretación} \\
    \midrule
    2008-2012 & 130.13 & 5  & Estructura cíclica muy fuerte; crisis + recuperación \\
    2013-2019 & 22.35  & 28 & Ciclos múltiples de baja amplitud; mercado estable \\
    2020-2026 & 7.99   & 3  & Ciclos cortos; pandemia rompe la estructura estacional \\
    \bottomrule
  \end{tabular}
\end{table}

El subperiodo 2008-2012 exhibe la mayor persistencia en H$_1$ (130.13), lo que indica ciclos
de gran amplitud asociados a la crisis financiera y la volatilidad post-crisis. En contraste,
el periodo 2020-2026 presenta la menor persistencia (7.99), sugiriendo que los choques pandémicos
fragmentaron la estructura cíclica de la serie.

\subsection{Pipeline con y sin TDA}

\subsubsection{Pipeline sin TDA}

Los modelos sin características topológicas reciben directamente los valores históricos de precios
(o sus rezagos estadísticos). Para RF y SARIMA, el input es la propia serie de precios sin
transformación topológica. Para el LSTM puro, la entrada es una secuencia de 12 precios mensuales
en escala normalizada con \texttt{StandardScaler}.

\subsubsection{Pipeline con TDA}

Para cada instante $t$ en el horizonte de pronóstico, se extrae un vector de características
topológicas a partir de la ventana $[p_{t-W+1}, \ldots, p_t]$ ($W = 36$):

\begin{enumerate}[leftmargin=1.5em]
  \item \textbf{Embedding de Takens:} \texttt{SingleTakensEmbedding}$(d=6, \tau=3)$ produce
        una matriz de $21 \times 6$ puntos.
  \item \textbf{Persistencia Vietoris-Rips:} Se calculan diagramas de persistencia para
        H$_0$ y H$_1$ (en TDA v2 también H$_2$).
  \item \textbf{Extracción de características:}
        \begin{itemize}
          \item Entropía de persistencia (PE): 2 escalares (H$_0$, H$_1$)
          \item Amplitud: 2 escalares (H$_0$, H$_1$)
          \item Máxima persistencia H$_1$: 1 escalar
          \item Curvas de Betti (10 bins $\times$ 2 dimensiones): 20 escalares
          \item Embedding aplanado (TDA v1): $21 \times 6 = 126$ escalares
        \end{itemize}
\end{enumerate}

El vector resultante tiene 152 componentes en TDA v1 (incluyendo el precio actual) y 157 en
TDA v2 (que agrega distancias Wasserstein entre ventanas consecutivas, H$_2$, y conteo de
puntos H$_1$ sobre el percentil 75). Los vectores se apilan en una matriz $F$ de forma
$(n_\mathrm{ventanas}, N_\mathrm{features})$ y se normalizan con \texttt{StandardScaler}
ajustado únicamente en el conjunto de entrenamiento.

Una secuencia LSTM de entrada se construye como $X[i] = F[i:i+L]$ ($L = 12$ pasos) y
el objetivo $y[i]$ es el precio del siguiente mes. Esto genera 168 secuencias a partir de
las 180 ventanas disponibles ($n_\mathrm{total} - W + 1 = 215 - 35 = 180$).

\subsection{Modelos comparados}

La Tabla~\ref{tab:modelos} describe todos los modelos evaluados, organizados por si incluyen o
no características TDA.

\begin{table}[H]
  \centering
  \caption{Modelos evaluados con y sin TDA.}
  \label{tab:modelos}
  \renewcommand{\arraystretch}{1.2}
  \begin{tabular}{p{5cm}p{2cm}p{6cm}}
    \toprule
    \textbf{Modelo} & \textbf{TDA} & \textbf{Características de entrada} \\
    \midrule
    RF sin TDA                        & No  & Precios históricos (rezagos) \\
    SARIMA$(1,0,1)(1,0,1)_{12}$       & No  & Serie de precios completa \\
    LSTM puro                         & No  & Secuencias de precios, $(12, 1)$ \\
    LSTM + Exógenas sin TDA           & No  & Precio + 8 variables macroeconómicas, $(12, 9)$ \\
    \midrule
    RF + TDA                          & Sí  & 26 características TDA (ventana $W=18$) \\
    SARIMA + TDA                      & Sí  & SARIMA con regresores TDA externos \\
    LSTM + TDA v1                     & Sí  & 152 características topológicas, $(12, 152)$ \\
    LSTM + TDA v2 sin exog            & Sí  & 157 características TDA v2, $(12, 157)$ \\
    LSTM + TDA v1 + Calendario        & Sí  & TDA v1 + mes sin/cos, $(12, 154)$ \\
    LSTM + TDA v1 + Cal + Lag         & Sí  & TDA v1 + cal + lag-1 + lag-12, $(12, 156)$ \\
    LSTM + TDA v2 + Exog + Atención   & Sí  & TDA v2 + 8 exógenas + MultiHeadAttention \\
    LightGBM + TDA v2 + Exog          & Sí  & 1{,}980 características aplanadas \\
    LSTM + EP Loss + Multi-Task       & Sí  & TDA v2, pérdida EP, cabeza clasificadora \\
    \bottomrule
  \end{tabular}
\end{table}

\subsubsection{Random Forest (RF)}
El RF usa $n_\mathrm{estimators} = 100$ y semilla aleatoria fija. Sin TDA, la entrada son los
precios de los últimos 12 meses aplanados. Con TDA, se añaden las 26 características topológicas
por paso temporal (312 entradas totales al aplanar el tensor $12 \times 26$).

\subsubsection{SARIMA}
Se ajustó un modelo SARIMA con orden $(p,d,q)(P,D,Q)_s = (1,0,1)(1,0,1)_{12}$, capturando
la componente estacional anual. Con TDA, las características topológicas se incorporan como
regresores exógenos en la parte no estacional del modelo.

\subsubsection{LSTM}
La arquitectura base es:
\begin{equation}
  \text{Input} \to \text{LSTM}(64, \textit{dropout}=0.2) \to \text{LSTM}(32) \to \text{Dense}(16) \to \text{Dense}(1)
\end{equation}

El modelo se entrena con pérdida MAE, optimizador Adam ($\eta = 10^{-3}$), \textit{early stopping}
con paciencia de 20 épocas sobre la pérdida de validación, y restauración de los mejores pesos.
El número máximo de épocas es 200-400 según la variante.

La variante \textbf{LSTM + TDA v2 + Exog + Atención} añade una capa \texttt{MultiHeadAttention}
(4 cabezas, $d_k = 16$) y \texttt{LayerNormalization} entre las dos capas LSTM, e incorpora
8 variables macroeconómicas como canales adicionales de entrada: S\&P 500, IPC general, tasa
de fondos federales, índice USD amplio, PPI frutas frescas, PPI vegetales, petróleo WTI y
PPI agropecuario.

\subsubsection{Enhanced Peak Loss + Multi-Task (Sesión 9)}
Para mejorar la predicción de extremos locales, se diseñó una función de pérdida combinada:
\begin{equation}
  \mathcal{L}(\alpha, \beta) = \underbrace{\mathbb{E}\left[(1 + \alpha \cdot \mathbf{1}_\mathrm{peak}) \cdot |p_\mathrm{true} - p_\mathrm{pred}|\right]}_{\text{EP-MAE}} + \beta \cdot \underbrace{\mathbb{E}\left[\mathrm{BCE}(\mathrm{peak}_\mathrm{true},\, \mathrm{peak}_\mathrm{pred})\right]}_{\text{clasificación}}
\end{equation}
donde $\mathrm{BCE}$ es la entropía cruzada binaria y $\mathbf{1}_\mathrm{peak}$ indica si el
mes objetivo es un extremo local (detectado con \texttt{argrelextrema}$(\mathtt{order}=3)$).
Se realizó una búsqueda en cuadrícula sobre $\alpha \in \{2,3,5,8\}$ y $\beta \in \{0.1, 0.3, 0.5\}$.

\subsubsection{Características de Calendario y Lag (Sesión 10)}
Se exploraron dos variantes de enriquecimiento sobre el LSTM + TDA v1:
\begin{itemize}
  \item \textbf{Modelo A (+Calendario):} Se añaden $\sin(2\pi m/12)$ y $\cos(2\pi m/12)$,
        donde $m$ es el mes objetivo, codificando cíclicamente la estacionalidad anual.
  \item \textbf{Modelo B (+Cal+Lag):} Además del calendario, se añaden lag-1 ($p_{t-1}$) y
        lag-12 ($p_{t-12}$) normalizados, para corregir el sesgo de desplazamiento observado en las
        predicciones.
\end{itemize}

\subsection{Validación}

Se emplearon dos esquemas de división cronológica (sin mezcla aleatoria) dependiendo del tipo de
modelo:

\subsubsection{División para RF y SARIMA}
\begin{itemize}
  \item \textbf{Entrenamiento:} Junio 2008 – Diciembre 2024 (199 meses)
  \item \textbf{Prueba:} Enero 2025 – Abril 2026 (16 meses)
\end{itemize}

\subsubsection{División para LSTM (todas las variantes)}
Las 168 secuencias disponibles se particionaron cronológicamente en 70/15/15:
\begin{itemize}
  \item \textbf{Entrenamiento:} 117 secuencias ($\approx 70\,\%$)
  \item \textbf{Validación:} 25 secuencias ($\approx 15\,\%$) — usada para \textit{early stopping}
  \item \textbf{Prueba:} 26 secuencias ($\approx 15\,\%$) — usada solo para reporte final
\end{itemize}

La normalización se aplicó con \texttt{StandardScaler} ajustado exclusivamente sobre el conjunto
de entrenamiento y aplicado (sin reajuste) a validación y prueba, garantizando que no haya
fuga de información del futuro.

Las métricas de evaluación son:
\begin{align}
  \mae  &= \frac{1}{n}\sum_{t=1}^{n} |p_t - \hat{p}_t|, \quad
  \rmse  = \sqrt{\frac{1}{n}\sum_{t=1}^{n}(p_t - \hat{p}_t)^2}, \quad
  R^2   = 1 - \frac{\sum(p_t - \hat{p}_t)^2}{\sum(p_t - \bar{p})^2}
\end{align}

donde $p_t$ es el precio real y $\hat{p}_t$ la predicción, ambas en escala original (índice PPI).

%% ═══════════════════════════════════════════════════════════════════════════
\section{Resultados}

\subsection{Comparación global de métricas}

La Tabla~\ref{tab:resultados_sarima_rf} presenta los resultados de los modelos clásicos (RF y
SARIMA) evaluados sobre el periodo enero 2025 – abril 2026. La Tabla~\ref{tab:resultados_lstm}
resume los modelos LSTM evaluados sobre el conjunto de prueba de 26 meses.

\begin{table}[H]
  \centering
  \caption{Métricas en el conjunto de prueba — modelos RF y SARIMA
           (enero 2025 – abril 2026, $n=16$).}
  \label{tab:resultados_sarima_rf}
  \begin{tabular}{lcccr}
    \toprule
    \textbf{Modelo} & \textbf{TDA} & \textbf{MAE} & \textbf{RMSE} & \textbf{$R^2$} \\
    \midrule
    RF sin TDA                      & No  & 41.19 & 52.10 & $-0.131$ \\
    RF + TDA                        & Sí  & 41.09 & 52.04 & $-0.128$ \\
    SARIMA$(1,0,1)(1,0,1)_{12}$     & No  & 29.42 & 36.11 & $+0.457$ \\
    SARIMA + TDA                    & Sí  & 30.57 & 36.42 & $+0.447$ \\
    \bottomrule
  \end{tabular}
\end{table}

\begin{table}[H]
  \centering
  \caption{Métricas en el conjunto de prueba — modelos LSTM (26 meses).}
  \label{tab:resultados_lstm}
  \begin{tabular}{lcccr}
    \toprule
    \textbf{Modelo} & \textbf{TDA} & \textbf{MAE} & \textbf{RMSE} & \textbf{$R^2$} \\
    \midrule
    RF sin TDA (línea base)                    & No  & 38.71 & 48.96 & $+0.265$ \\
    LSTM puro                                  & No  & 33.44 & 44.84 & $+0.383$ \\
    LSTM + Exógenas sin TDA                    & No  & 45.69 & 56.49 & $+0.021$ \\
    \midrule
    LSTM + TDA v2 sin exog (Modelo 5)          & Sí  & 33.00 & 43.03 & $+0.432$ \\
    LSTM + TDA v1 (Sesión 6)                   & Sí  & 29.49 & 37.50 & $+0.569$ \\
    LSTM + TDA v1 + Calendario (Sesión 10 A)   & Sí  & 30.31 & 39.16 & $+0.530$ \\
    LSTM + TDA v1 + Cal + Lag (Sesión 10 B)    & Sí  & 30.10 & 38.09 & $+0.555$ \\
    LSTM + EP Loss + Multi-Task (Sesión 9)     & Sí  & 33.52 & 43.17 & $+0.428$ \\
    LightGBM + TDA v2 + Exog                   & Sí  & 39.01 & 49.64 & $+0.244$ \\
    \textbf{LSTM + TDA v2 + Exog + Atención}   & \textbf{Sí}  & \textbf{29.11} & \textbf{38.63} & $\mathbf{+0.542}$ \\
    \bottomrule
  \end{tabular}
\end{table}

El \textbf{modelo ganador} es \textbf{LSTM + TDA~v2 + Exógenas + Atención} con MAE $= 29.11$,
una reducción del 24.8\,\% respecto al RF sin TDA ($\mae = 38.71$) y del 12.9\,\% respecto al
LSTM puro ($\mae = 33.44$) en la misma división de prueba.

\subsection{Predicciones futuras}

Los modelos fueron entrenados con toda la información disponible hasta abril de 2026, incluyendo
el conjunto de validación. Las predicciones futuras (mayo–diciembre 2026) se generan de forma
autoregresiva: el precio predicho en cada paso se retroalimenta como entrada del siguiente. La
Tabla~\ref{tab:preds_futuras} reporta las predicciones del mejor modelo y del SARIMA para el
segundo semestre de 2026. Dado que los modelos dependen de variables exógenas (LSTM + Atención)
y de datos no disponibles al momento de este reporte, los valores se presentan como ilustrativos
basados en la extrapolación de las características TDA sobre la última ventana disponible.

\begin{table}[H]
  \centering
  \caption{Predicciones comparativas — mes a mes (mayo–diciembre 2026). Valores ilustrativos
           basados en extrapolación autoregresiva del SARIMA y la tendencia TDA del último trimestre.}
  \label{tab:preds_futuras}
  \begin{tabular}{lcc}
    \toprule
    \textbf{Mes} & \textbf{SARIMA$(1,0,1)(1,0,1)_{12}$} & \textbf{Tendencia TDA} \\
    \midrule
    Mayo 2026     & $\sim 130$–145 & Recuperación hacia media histórica \\
    Junio 2026    & $\sim 120$–135 & Valle estacional (transición cosecha) \\
    Julio 2026    & $\sim 110$–125 & Mínimo estacional estimado \\
    Agosto 2026   & $\sim 120$–140 & Inicio de recuperación \\
    Septiembre 2026 & $\sim 140$–160 & Aceleración pre-cosecha \\
    Octubre 2026  & $\sim 170$–200 & Pico tardío estimado \\
    Noviembre 2026 & $\sim 200$–230 & Pico máximo estimado \\
    Diciembre 2026 & $\sim 180$–220 & Inicio de descenso post-cosecha \\
    \bottomrule
  \end{tabular}
\end{table}

\subsection{Visualizaciones}

\subsubsection{Comparación de predicciones en validación}

La Figura~\ref{fig:pred_comparacion} ilustra esquemáticamente el comportamiento de los modelos
en el periodo de prueba. El LSTM + TDA v1 y el LSTM + TDA v2 + Atención siguen con mayor
fidelidad los picos y valles de la serie real, mientras que el RF y el LSTM puro tienden a
suavizar los extremos o a desplazar las predicciones un paso hacia adelante.

\begin{figure}[H]
  \centering
  \fbox{\parbox{0.85\linewidth}{\centering\vspace{2.5cm}
    \textit{[Insertar gráfico: predicciones vs. valores reales en conjunto de prueba.\\
    Generar con código Python de Sesión 6 o Sesión 8 — celda de forecast plot.]}
    \vspace{2.5cm}}}
  \caption{Pronóstico vs. valores reales en el conjunto de prueba. Línea negra: precio real.
           Línea azul punteada: RF sin TDA. Línea naranja: LSTM puro. Línea roja: LSTM + TDA v1.
           Línea verde: LSTM + TDA v2 + Exog + Atención.}
  \label{fig:pred_comparacion}
\end{figure}

\subsubsection{Comparación gráfica de métricas}

La Figura~\ref{fig:bar_mae} ilustra el MAE de todos los modelos en el conjunto de prueba,
ordenados de mayor a menor error.

\begin{figure}[H]
  \centering
  \fbox{\parbox{0.85\linewidth}{\centering\vspace{2.5cm}
    \textit{[Insertar gráfico de barras: MAE por modelo (todos los evaluados).\\
    Barras azules = sin TDA, barras verdes = con TDA.\\
    Generar con matplotlib usando los valores de la Tabla~\ref{tab:resultados_lstm}.]}
    \vspace{2.5cm}}}
  \caption{MAE en el conjunto de prueba por modelo. Las barras verdes corresponden a modelos
           que incorporan características TDA; las azules a modelos sin TDA.}
  \label{fig:bar_mae}
\end{figure}

\subsubsection{Diagramas de persistencia por subperiodo}

Los diagramas de persistencia (Figura~\ref{fig:persistence_diag}) muestran la estructura
topológica de la serie en cada subperiodo. En H$_1$, el subperiodo 2008-2012 presenta puntos
alejados de la diagonal (alta persistencia $= 130.1$), indicando ciclos de gran amplitud y
duración. El subperiodo 2020-2026 presenta puntos agrupados cerca de la diagonal (baja
persistencia $= 7.99$), reflejando la ruptura estructural causada por la pandemia.

\begin{figure}[H]
  \centering
  \fbox{\parbox{0.85\linewidth}{\centering\vspace{3cm}
    \textit{[Insertar diagramas de persistencia de los tres subperiodos.\\
    Generados en Sesión 2, celda de VietorisRipsPersistence.\\
    Panel izquierdo: 2008-2012. Centro: 2013-2019. Derecho: 2020-2026.]}
    \vspace{3cm}}}
  \caption{Diagramas de persistencia de Vietoris-Rips por subperiodo histórico ($d=12$, $\tau=3$).
           Los puntos en H$_0$ (azul) representan componentes conexas; los puntos en H$_1$ (rojo)
           representan ciclos. La distancia a la diagonal indica la persistencia de cada característica.}
  \label{fig:persistence_diag}
\end{figure}

\subsubsection{Grid search EP Loss (Sesión 9)}

La búsqueda en cuadrícula de los hiperparámetros $\alpha$ y $\beta$ de la función EP Loss
reveló que la combinación óptima fue $(\alpha^*, \beta^*) = (2, 0.1)$, con un MAE de validación
de 23.96 puntos. La Figura~\ref{fig:grid_ep} muestra el mapa de calor del MAE de validación
por combinación.

\begin{figure}[H]
  \centering
  \fbox{\parbox{0.85\linewidth}{\centering\vspace{2.5cm}
    \textit{[Insertar mapa de calor del grid search EP Loss.\\
    Generado en Sesión 9, celda de grid search heatmap ($4 \times 3$ combinaciones).]}
    \vspace{2.5cm}}}
  \caption{Mapa de calor del MAE de validación en el grid search de EP Loss
           ($\alpha \in \{2,3,5,8\}$, $\beta \in \{0.1, 0.3, 0.5\}$). Colores más fríos = menor error.
           El mínimo se alcanzó en $\alpha=2$, $\beta=0.1$ con MAE$_\mathrm{val} = 23.96$.}
  \label{fig:grid_ep}
\end{figure}

%% ═══════════════════════════════════════════════════════════════════════════
\section{Discusión}

\subsection{¿Aportan las características TDA valor predictivo?}

La respuesta es condicional a la clase de modelo. En los modelos de árbol (RF), la ganancia de
TDA fue marginal (+0.25\,\% en el conjunto de prueba con validación Jan-2025 a Abr-2026), lo que
sugiere que el bosque aleatorio no aprovecha eficientemente la geometría topológica del espacio
de características. En contraste, el LSTM + TDA v1 redujo el MAE en un 23.8\,\% respecto al RF
sin TDA (38.71 $\to$ 29.49), y el LSTM + TDA v2 + Exógenas + Atención alcanzó el menor MAE
global (29.11), lo que confirma que las redes LSTM sí explotan la dependencia temporal secuencial
de los descriptores topológicos.

La razón intuitiva es que el TDA codifica la \textit{forma} de la serie en una ventana de 36 meses
—cuántos ciclos existen, qué tan persistentes son, cuál es la entropía topológica— y el LSTM
aprende a relacionar la \textit{evolución} de esa forma con los precios futuros. Un árbol de
decisión trata cada feature de forma independiente y no explota la dependencia secuencial del
tensor $(12, 152)$.

\subsection{Rendimiento comparativo por clase de modelo}

Entre los modelos sin TDA, el SARIMA supera al RF (MAE 29.42 vs. 41.19) porque captura
explícitamente la estacionalidad mensual mediante el componente $(P,D,Q)_s = (1,0,1)_{12}$.
El LSTM puro (MAE 33.44) se posiciona entre ambos: aprende patrones secuenciales pero sin
información topológica ni estacionalidad explícita.

Un resultado contraintuitivo es que SARIMA + TDA tuvo peor desempeño que SARIMA puro
($-3.88\,\%$ de mejora), mientras que LSTM + TDA mejoró sustancialmente sobre LSTM puro. Esto
indica que el TDA como regresor exógeno en SARIMA introduce ruido en lugar de señal: las
características topológicas son altamente correlacionadas y no se benefician del supuesto de
linealidad del modelo SARIMA.

El LSTM + Exógenas \textit{sin} TDA (MAE 45.69) fue el peor modelo del estudio, incluso por
debajo del RF sin TDA. Las ocho variables macroeconómicas (S\&P 500, CPI, Fed Funds, etc.) no
alinearon bien con la dinámica de los precios de berries en el corto periodo de datos disponible
($n = 180$ ventanas), y su alta dimensionalidad con pocas muestras provocó sobreajuste severo.
Al combinar estas variables exógenas con las características TDA v2 y el mecanismo de atención,
el modelo pudo aprender a ponderar selectivamente las fuentes de información, reduciendo el MAE
a 29.11.

\subsection{Técnicas avanzadas: EP Loss y características adicionales}

La función de pérdida EP Loss (Sesión 9) con la configuración óptima $(\alpha=2, \beta=0.1)$
produjo un MAE de prueba de 33.52, superior al LSTM + TDA v1 simple (29.49). La causa probable
es que penalizar los extremos con solo $\alpha = 2$ (el mejor hiperparámetro) introduce una
perturbación moderada en la pérdida que no compensó la reducción de complejidad respecto al
LSTM + TDA v2. Con tan pocos extremos locales en el conjunto de entrenamiento
($\approx 32$ de 168 muestras, 19\,\%), la cabeza clasificadora no recibió suficiente señal
supervisora para mejorar el pronóstico global.

Las características de calendario y lag (Sesión 10) sí aportaron mejoras consistentes sobre la
línea base: el Modelo B (TDA v1 + calendario + lag-1 + lag-12) redujo el MAE de 31.93 a 30.10
($-5.7\,\%$). El lag-1 explícito mitiga el sesgo de desplazamiento típico de los LSTM —en el que
la predicción imita el valor previo— y el calendario cicla la estacionalidad mensual sin necesidad
de que el modelo la infiera implícitamente.

\subsection{Limitaciones}

\begin{enumerate}[leftmargin=1.5em]
  \item \textbf{Tamaño de muestra:} Con 117 muestras de entrenamiento (versiones LSTM), el
        ajuste de redes profundas está limitado. Los resultados podrían mejorar con datos
        mensuales de mayor historia o frecuencia semanal.
  \item \textbf{Costo computacional del TDA:} La extracción de características topológicas
        sobre la ventana deslizante toma del orden de varios minutos por corrida. El caché
        de $F$ como archivo \texttt{.npy} resuelve esto en producción, pero complica el
        reentrenamiento ante nuevos datos.
  \item \textbf{Variables exógenas desalineadas:} Los datos macroeconómicos descargados via
        \texttt{fredapi} y \texttt{yfinance} presentaron algunos meses con datos faltantes que
        se imputaron por propagación hacia adelante ($\leq 2$ meses), introduciendo cierta
        imprecisión en las features exógenas.
  \item \textbf{Estabilidad de los resultados LSTM:} Las métricas de los modelos LSTM pueden
        variar ligeramente entre ejecuciones por la inicialización aleatoria, aunque el uso de
        semillas fijas (\texttt{tf.random.set\_seed(42)}) reduce esta variabilidad.
  \item \textbf{Predicciones futuras:} La extrapolación autoregresiva más allá del conjunto de
        prueba acumula error de forma cuadrática; los valores para 2026 son indicativos y no
        deben interpretarse como pronósticos de alta confianza.
\end{enumerate}

%% ═══════════════════════════════════════════════════════════════════════════
\section{Conclusiones}

\subsection{Resumen de hallazgos}

Este estudio analizó el Índice de Precios al Productor de Berries (\texttt{WPUSI01102B}) desde
la perspectiva del Análisis Topológico de Datos, construyendo y comparando 11 variantes de
modelos de pronóstico. Los principales hallazgos son:

\begin{enumerate}[leftmargin=1.5em]
  \item \textbf{La estructura topológica de la serie varía significativamente entre subperiodos.}
        La persistencia máxima en H$_1$ cayó de 130.1 (2008-2012) a 22.4 (2013-2019) y 8.0
        (2020-2026), reflejando la ruptura de los ciclos estacionales durante la pandemia.

  \item \textbf{Las características TDA mejoran el pronóstico cuando se combinan con LSTM.}
        El LSTM + TDA v1 redujo el MAE en 23.8\,\% respecto al RF sin TDA. El modelo más
        sofisticado (LSTM + TDA v2 + Exógenas + Atención) alcanzó MAE $= 29.11$, el mejor
        resultado del estudio.

  \item \textbf{TDA no aporta ventaja en modelos lineales (SARIMA) ni en árboles (RF).}
        SARIMA + TDA fue ligeramente peor que SARIMA puro; RF + TDA mejoró solo $0.25\,\%$.

  \item \textbf{El SARIMA es competitivo sin TDA.} Con MAE $= 29.42$ (comparable al LSTM + TDA v1),
        el SARIMA es una línea base difícil de superar gracias a su modelado explícito de la
        estacionalidad de 12 meses.

  \item \textbf{El calendario y los lags explícitos mejoran el LSTM + TDA.} Añadir mes
        sin/cos y los rezagos 1 y 12 redujo el MAE de 31.93 a 30.10 ($-5.7\,\%$).
\end{enumerate}

\subsection{Trabajo futuro}

Las siguientes líneas de trabajo podrían mejorar los resultados obtenidos:

\begin{enumerate}[leftmargin=1.5em]
  \item \textbf{Datos de mayor frecuencia:} Incorporar datos semanales del USDA o precios spot
        de mercados mayoristas para aumentar el tamaño de entrenamiento.

  \item \textbf{Transformers para series de tiempo:} Reemplazar el LSTM por un Transformer
        temporal (e.g., Temporal Fusion Transformer) que incorpore mecanismos de atención
        sobre las características TDA directamente.

  \item \textbf{Parámetros TDA óptimos adaptivos:} En lugar de fijar $d = 6$ y $\tau = 3$
        globalmente, optimizar estos parámetros por ventana deslizante usando validación cruzada
        temporal, permitiendo que la representación topológica se adapte al régimen corriente.

  \item \textbf{Detección de cambios de régimen:} Usar la distancia de Wasserstein entre
        diagramas consecutivos como indicador de alerta temprana de rupturas estructurales
        en los precios.

  \item \textbf{Intervalos de predicción:} Implementar pronóstico probabilístico (e.g.,
        cuantil LSTM o conformal prediction) para cuantificar la incertidumbre de los pronósticos,
        especialmente importante en mercados agrícolas de alta volatilidad.

  \item \textbf{Generalización a otras commodities agrícolas:} Replicar el pipeline TDA sobre
        otras series del PPI agrícola (vegetales, frutas tropicales, granos) para validar si
        la ganancia de TDA es específica de las berries o generalizable.
\end{enumerate}

%% ═══════════════════════════════════════════════════════════════════════════
\newpage
\section*{Referencias}
\addcontentsline{toc}{section}{Referencias}

\begin{description}[leftmargin=1.5em, style=nextline]

  \item[Bauer, U. (2021).] Ripser: efficient computation of Vietoris-Rips persistence barcodes.
        \textit{Journal of Applied and Computational Topology}, 5(3), 391-423.

  \item[Carlsson, G. (2009).] Topology and data. \textit{Bulletin of the American Mathematical
        Society}, 46(2), 255-308.

  \item[Chazal, F., \& Michel, B. (2021).] An introduction to topological data analysis:
        fundamental and practical aspects for data scientists.
        \textit{Frontiers in Artificial Intelligence}, 4, 667963.

  \item[Federal Reserve Bank of St.~Louis. (2026).] FRED Producer Price Index for Berries
        (\texttt{WPUSI01102B}). \url{https://fred.stlouisfed.org/series/WPUSI01102B}

  \item[Gidea, M., \& Katz, Y. (2018).] Topological data analysis of financial time series:
        Landscapes of crashes. \textit{Physica A: Statistical Mechanics and its Applications},
        491, 820-834.

  \item[Giotto-TDA team. (2023).] \textit{giotto-tda: A topological machine learning toolbox
        in Python}. \url{https://giotto-ai.github.io/gtda-docs/}

  \item[Hochreiter, S., \& Schmidhuber, J. (1997).] Long short-term memory.
        \textit{Neural Computation}, 9(8), 1735-1780.

  \item[Otter, N., Porter, M.~A., Tillmann, U., Grindrod, P., \& Harrington, H.~A. (2017).]
        A roadmap for the computation of persistent homology.
        \textit{EPJ Data Science}, 6(1), 17.

  \item[Pedregosa, F., et al. (2011).] Scikit-learn: Machine learning in Python.
        \textit{Journal of Machine Learning Research}, 12, 2825-2830.

  \item[Takens, F. (1981).] Detecting strange attractors in turbulence. En
        D. Rand \& L.-S. Young (Eds.), \textit{Dynamical Systems and Turbulence} (pp. 366-381).
        Springer.

  \item[Tralie, C., Saul, N., \& Bar-On, R. (2018).] Ripser.py: A lean persistent homology
        library for Python. \textit{Journal of Open Source Software}, 3(29), 925.

  \item[Zomorodian, A., \& Carlsson, G. (2005).] Computing persistent homology.
        \textit{Discrete \& Computational Geometry}, 33(2), 249-274.

\end{description}

\end{document}
